\documentclass[12pt]{article}
\usepackage[utf8]{inputenc}
\usepackage[T1]{fontenc}
\usepackage{hyperref}
\title{Zapotec, Xadani Language Resources for K-12 Curriculum Development}
\author{Compiled by Claude Sonnet 5.5 (claude-sonnet-5-5) for Andr\'as Kornai}
\date{}
\begin{document}
\maketitle

\section{Introduction}
Xadani Zapotec (ISO 639-3 zax, Glottocode xada1235), also called Eastern Pochutla Zapotec or Zapoteco de Santa Mar\'ia Xadani, is an Otomanguean language of Oaxaca, Mexico. Joshua Project estimates about 500 speakers. Very little exists, and no dictionary, grammar or school material was found in the search results.

\section{Primary Educational Resources}
None found.

\section{Linguistic and Reference Materials}
No dictionary or grammar was found. Wikitongues lists the language as unwritten, with an EGIDS status of 7 (Shifting).
\begin{itemize}
\item Wikitongues page: \url{https://wikitongues.org/languages/zax/}
\end{itemize}

\section{Religious and Community Resources}
Joshua Project reports no Bible portions, no New Testament, no Jesus Film and no audio Bible, and states that no resources are reported in this language.
\begin{itemize}
\item Joshua Project: \url{https://joshuaproject.net/languages/zax}
\end{itemize}

\section{Digital and Audio Resources}
Global Recordings Network has no recordings in Xadani Zapotec.
\begin{itemize}
\item Global Recordings Network: \url{https://globalrecordings.net/en/language/zax}
\end{itemize}

\section{Curriculum Development Context}
The language is endangered and is used as a first language by adults only. It is not known to be taught in schools and has no established written tradition. A curriculum would need new documentation done with the community, starting with an orthography and a recorded corpus.

\end{document}